% !TEX root = main.tex

Pre-trained convolutional neural networks, or convnets, have become the
building blocks in most computer vision
applications~\cite{ren2015faster,chen2016deeplab,weinzaepfel2013deepflow,carreira2016human}.
They produce excellent general-purpose features that can be used to improve the
generalization of models learned on a limited amount of
data~\cite{sharif2014cnn}.  The existence of ImageNet~\cite{deng2009imagenet},
a large fully-supervised dataset, has been fueling advances in pre-training of
convnets.  However, Stock and Cisse~\cite{stock2017convnets} have recently
presented empirical evidence that the performance of state-of-the-art classifiers on
ImageNet is largely underestimated, and little error is left unresolved.  This
explains in part why the performance has been saturating despite the numerous novel
architectures proposed in recent
years~\cite{chen2016deeplab,he2015delving,huang2016densely}.  As a matter of
fact, ImageNet is relatively small by today's standards; it ``only'' contains a
million images that cover the specific domain of object classification.  A
natural way to move forward is to build a bigger and more diverse dataset,
potentially consisting of billions of images.  This, in turn, would require a
tremendous amount of manual annotations, despite the expert knowledge in
crowdsourcing accumulated  by the community over the
years~\cite{kovashka2016crowdsourcing}.  Replacing labels by raw metadata
leads to biases in the visual representations with unpredictable
consequences~\cite{misra2016seeing}.  This calls for methods that can be
trained on internet-scale datasets with no supervision.

Unsupervised learning has been widely studied in the Machine Learning
community~\cite{friedman2001elements}, and algorithms for clustering,
dimensionality reduction or density estimation are regularly used in computer
vision
applications~\cite{turk1991face,shi2000normalized,joulin2010discriminative}.
For example, the ``bag of features'' model uses clustering on handcrafted local
descriptors to produce good image-level features~\cite{csurka2004visual}.  A
key reason for their success is that they can be applied on any specific domain
or dataset, like satellite or medical images, or on images captured with a new
modality, like depth, where annotations are not always available in quantity.
Several works have shown that it was possible to adapt unsupervised methods
based on density estimation or dimensionality reduction to deep
models~\cite{goodfellow2014generative,kingma2013auto}, leading to promising
all-purpose visual
features~\cite{bojanowski2017unsupervised,donahue2016adversarial}.  Despite the
primeval success of clustering approaches in image classification, very few
works~\cite{yang2016joint,xie2016unsupervised} have been proposed to adapt them
to the end-to-end training of convnets, and never at scale.  An issue is that
clustering methods have been primarily designed for linear models on top of
fixed features, and they scarcely work if the features have to be learned
simultaneously.  For example, learning a convnet with $k$-means would lead to a
trivial solution where the features are zeroed, and the clusters are collapsed
into a single entity.

\begin{figure}[!t]
  \centering
  \includegraphics[width=.9\linewidth]{images-pullfig4bis_clean.jpg}
  \caption{
    Illustration of the proposed method: we iteratively cluster deep features and use the cluster assignments as pseudo-labels to learn the parameters of the convnet.
  }
  \label{fig:pullfig}
\end{figure}

In this work, we propose a novel clustering approach for the large scale
end-to-end training of convnets.  We show that it is possible to obtain useful
general-purpose visual features with a clustering framework.  Our approach,
summarized in Figure~\ref{fig:pullfig}, consists in alternating between
clustering of the image descriptors and updating the weights of the convnet by
predicting the cluster assignments.  For simplicity, we focus our study on
$k$-means, but other clustering approaches can be used, like Power Iteration
Clustering~(PIC)~\cite{lin2010power}.  The overall pipeline is sufficiently
close to the standard supervised training of a convnet to reuse many common
tricks~\cite{ioffe2015batch}.  Unlike self-supervised
methods~\cite{doersch2015unsupervised,noroozi2016unsupervised,pathak2016learning},
clustering has the advantage of requiring little domain knowledge and no specific
signal from the
inputs~\cite{zhang2016colorful,wang2015unsupervised}.  Despite its simplicity,
our approach achieves significantly higher performance than previously
published unsupervised methods on both ImageNet classification and transfer tasks.

Finally, we probe the robustness of our framework by modifying the experimental protocol,
in particular the training set and the convnet architecture. The resulting set of experiments extends the
discussion initiated by Doersch~\etal~\cite{doersch2015unsupervised} on the impact of these choices on
the performance of unsupervised methods.
We demonstrate that our
approach is robust to a change of architecture.
Replacing an AlexNet by a
VGG~\cite{simonyan2014very} significantly improves the quality of the features
and their subsequent transfer performance.  More importantly, we discuss
the use of ImageNet as a training set for unsupervised models.  While it helps
understanding the impact of the labels on the performance of a network,
ImageNet has a particular image distribution inherited from its use for a
fine-grained image classification challenge: it is composed of well-balanced classes and contains
a wide variety of dog breeds for example.  We consider, as an alternative, random Flickr
images from the YFCC100M dataset of Thomee~\etal~\cite{thomee2015new}.  We show
that our approach maintains state-of-the-art performance when trained on this
uncured data distribution.  Finally,  current benchmarks focus on
the capability of unsupervised convnets to capture class-level
information.  We propose to also evaluate them on image retrieval
benchmarks to measure their capability to capture instance-level
information.

In this paper, we make the following contributions:
\textbf{(i)} a novel unsupervised method for the end-to-end learning of convnets that works with any standard clustering algorithm, like $k$-means, and requires minimal additional steps;
\textbf{(ii)} state-of-the-art performance on many standard transfer tasks used in unsupervised learning;
\textbf{(iii)} performance above the previous state of the art when trained on an uncured image distribution;
\textbf{(iv)} a discussion about the current evaluation protocol in unsupervised feature learning.

